%!TeX root=../sensetop.tex
\chapter[Chapter \thechapter]{}
\lettrine[lraise=0.3]{A}{fter} a proper resistance on the part of Mrs Ferrars, just so violent and so steady as to preserve her from that reproach which she always seemed fearful of incurring, the reproach of being too amiable, Edward was admitted to her presence, and pronounced to be again her son.

Her family had of late been exceedingly fluctuating. For many years of her life she had had two sons; but the crime and annihilation of Edward a few weeks ago, had robbed her of one; the similar annihilation of Robert had left her for a fortnight without any; and now, by the resuscitation of Edward, she had one again.

In spite of his being allowed once more to live, however, he did not feel the continuance of his existence secure, till he had revealed his present engagement; for the publication of that circumstance, he feared, might give a sudden turn to his constitution, and carry him off as rapidly as before. With apprehensive caution therefore it was revealed, and he was listened to with unexpected calmness. Mrs Ferrars at first reasonably endeavoured to dissuade him from marrying Miss Dashwood, by every argument in her power;—told him, that in Miss Morton he would have a woman of higher rank and larger fortune;—and enforced the assertion, by observing that Miss Morton was the daughter of a nobleman with thirty thousand pounds, while Miss Dashwood was only the daughter of a private gentleman with no more than \textit{three}; but when she found that, though perfectly admitting the truth of her representation, he was by no means inclined to be guided by it, she judged it wisest, from the experience of the past, to submit—and therefore, after such an ungracious delay as she owed to her own dignity, and as served to prevent every suspicion of good-will, she issued her decree of consent to the marriage of Edward and Elinor.

What she would engage to do towards augmenting their income was next to be considered; and here it plainly appeared, that though Edward was now her only son, he was by no means her eldest; for while Robert was inevitably endowed with a thousand pounds a-year, not the smallest objection was made against Edward's taking orders for the sake of two hundred and fifty at the utmost; nor was anything promised either for the present or in future, beyond the ten thousand pounds, which had been given with Fanny.

It was as much, however, as was desired, and more than was expected, by Edward and Elinor; and Mrs Ferrars herself, by her shuffling excuses, seemed the only person surprised at her not giving more.

With an income quite sufficient to their wants thus secured to them, they had nothing to wait for after Edward was in possession of the living, but the readiness of the house, to which Colonel Brandon, with an eager desire for the accommodation of Elinor, was making considerable improvements; and after waiting some time for their completion, after experiencing, as usual, a thousand disappointments and delays from the unaccountable dilatoriness of the workmen, Elinor, as usual, broke through the first positive resolution of not marrying till every thing was ready, and the ceremony took place in Barton church early in the autumn.

The first month after their marriage was spent with their friend at the Mansion-house; from whence they could superintend the progress of the Parsonage, and direct every thing as they liked on the spot;—could chuse papers, project shrubberies, and invent a sweep. Mrs Jennings's prophecies, though rather jumbled together, were chiefly fulfiled; for she was able to visit Edward and his wife in their Parsonage by Michaelmas, and she found in Elinor and her husband, as she really believed, one of the happiest couples in the world. They had in fact nothing to wish for, but the marriage of Colonel Brandon and Marianne, and rather better pasturage for their cows.

They were visited on their first settling by almost all their relations and friends. Mrs Ferrars came to inspect the happiness which she was almost ashamed of having authorised; and even the Dashwoods were at the expense of a journey from Sussex to do them honour.

<I will not say that I am disappointed, my dear sister,> said John, as they were walking together one morning before the gates of Delaford House, <\textit{that} would be saying too much, for certainly you have been one of the most fortunate young women in the world, as it is. But, I confess, it would give me great pleasure to call Colonel Brandon brother. His property here, his place, his house, every thing is in such respectable and excellent condition! And his woods,—I have not seen such timber any where in Dorsetshire, as there is now standing in Delaford Hanger! And though, perhaps, Marianne may not seem exactly the person to attract him, yet I think it would altogether be advisable for you to have them now frequently staying with you, for as Colonel Brandon seems a great deal at home, nobody can tell what may happen; for, when people are much thrown together, and see little of anybody else,—and it will always be in your power to set her off to advantage, and so forth. In short, you may as well give her a chance: you understand me.>

\begin{figure}[tbph]
\centering
\includegraphics[width=\linewidth]{50respectable}
\caption{Every thing is in such respectable condition}
\end{figure}

But though Mrs Ferrars \textit{did} come to see them, and always treated them with the make-believe of decent affection, they were never insulted by her real favour and preference. \textit{That} was due to the folly of Robert, and the cunning of his wife; and it was earned by them before many months had passed away. The selfish sagacity of the latter, which had at first drawn Robert into the scrape, was the principal instrument of his deliverance from it; for her respectful humility, assiduous attentions, and endless flatteries, as soon as the smallest opening was given for their exercise, reconciled Mrs Ferrars to his choice, and re-established him completely in her favour.

The whole of Lucy's behaviour in the affair, and the prosperity which crowned it, therefore, may be held forth as a most encouraging instance of what an earnest, an unceasing attention to self-interest, however its progress may be apparently obstructed, will do in securing every advantage of fortune, with no other sacrifice than that of time and conscience. When Robert first sought her acquaintance, and privately visited her in Bartlett's Buildings, it was only with the view imputed to him by his brother. He merely meant to persuade her to give up the engagement; and as there could be nothing to overcome but the affection of both, he naturally expected that one or two interviews would settle the matter. In that point, however, and that only, he erred; for though Lucy soon gave him hopes that his eloquence would convince her in \textit{time}, another visit, another conversation, was always wanted to produce this conviction. Some doubts always lingered in her mind when they parted, which could only be removed by another half hour's discourse with himself. His attendance was by this means secured, and the rest followed in course. Instead of talking of Edward, they came gradually to talk only of Robert,—a subject on which he had always more to say than on any other, and in which she soon betrayed an interest even equal to his own; and in short, it became speedily evident to both, that he had entirely supplanted his brother. He was proud of his conquest, proud of tricking Edward, and very proud of marrying privately without his mother's consent. What immediately followed is known. They passed some months in great happiness at Dawlish; for she had many relations and old acquaintances to cut—and he drew several plans for magnificent cottages;—and from thence returning to town, procured the forgiveness of Mrs Ferrars, by the simple expedient of asking it, which, at Lucy's instigation, was adopted. The forgiveness, at first, indeed, as was reasonable, comprehended only Robert; and Lucy, who had owed his mother no duty and therefore could have transgressed none, still remained some weeks longer unpardoned. But perseverance in humility of conduct and messages, in self-condemnation for Robert's offence, and gratitude for the unkindness she was treated with, procured her in time the haughty notice which overcame her by its graciousness, and led soon afterwards, by rapid degrees, to the highest state of affection and influence. Lucy became as necessary to Mrs Ferrars, as either Robert or Fanny; and while Edward was never cordially forgiven for having once intended to marry her, and Elinor, though superior to her in fortune and birth, was spoken of as an intruder, \textit{she} was in every thing considered, and always openly acknowledged, to be a favourite child. They settled in town, received very liberal assistance from Mrs Ferrars, were on the best terms imaginable with the Dashwoods; and setting aside the jealousies and ill-will continually subsisting between Fanny and Lucy, in which their husbands of course took a part, as well as the frequent domestic disagreements between Robert and Lucy themselves, nothing could exceed the harmony in which they all lived together.

What Edward had done to forfeit the right of eldest son, might have puzzled many people to find out; and what Robert had done to succeed to it, might have puzzled them still more. It was an arrangement, however, justified in its effects, if not in its cause; for nothing ever appeared in Robert's style of living or of talking to give a suspicion of his regretting the extent of his income, as either leaving his brother too little, or bringing himself too much;—and if Edward might be judged from the ready discharge of his duties in every particular, from an increasing attachment to his wife and his home, and from the regular cheerfulness of his spirits, he might be supposed no less contented with his lot, no less free from every wish of an exchange.

Elinor's marriage divided her as little from her family as could well be contrived, without rendering the cottage at Barton entirely useless, for her mother and sisters spent much more than half their time with her. Mrs Dashwood was acting on motives of policy as well as pleasure in the frequency of her visits at Delaford; for her wish of bringing Marianne and Colonel Brandon together was hardly less earnest, though rather more liberal than what John had expressed. It was now her darling object. Precious as was the company of her daughter to her, she desired nothing so much as to give up its constant enjoyment to her valued friend; and to see Marianne settled at the mansion-house was equally the wish of Edward and Elinor. They each felt his sorrows, and their own obligations, and Marianne, by general consent, was to be the reward of all.

With such a confederacy against her—with a knowledge so intimate of his goodness—with a conviction of his fond attachment to herself, which at last, though long after it was observable to everybody else—burst on her—what could she do?

Marianne Dashwood was born to an extraordinary fate. She was born to discover the falsehood of her own opinions, and to counteract, by her conduct, her most favourite maxims. She was born to overcome an affection formed so late in life as at seventeen, and with no sentiment superior to strong esteem and lively friendship, voluntarily to give her hand to another!—and \textit{that} other, a man who had suffered no less than herself under the event of a former attachment, whom, two years before, she had considered too old to be married,—and who still sought the constitutional safeguard of a flannel waistcoat!

But so it was. Instead of falling a sacrifice to an irresistible passion, as once she had fondly flattered herself with expecting,—instead of remaining even for ever with her mother, and finding her only pleasures in retirement and study, as afterwards in her more calm and sober judgment she had determined on,—she found herself at nineteen, submitting to new attachments, entering on new duties, placed in a new home, a wife, the mistress of a family, and the patroness of a village.

Colonel Brandon was now as happy, as all those who best loved him, believed he deserved to be;—in Marianne he was consoled for every past affliction;—her regard and her society restored his mind to animation, and his spirits to cheerfulness; and that Marianne found her own happiness in forming his, was equally the persuasion and delight of each observing friend. Marianne could never love by halves; and her whole heart became, in time, as much devoted to her husband, as it had once been to Willoughby.

Willoughby could not hear of her marriage without a pang; and his punishment was soon afterwards complete in the voluntary forgiveness of Mrs Smith, who, by stating his marriage with a woman of character, as the source of her clemency, gave him reason for believing that had he behaved with honour towards Marianne, he might at once have been happy and rich. That his repentance of misconduct, which thus brought its own punishment, was sincere, need not be doubted;—nor that he long thought of Colonel Brandon with envy, and of Marianne with regret. But that he was for ever inconsolable, that he fled from society, or contracted an habitual gloom of temper, or died of a broken heart, must not be depended on—for he did neither. He lived to exert, and frequently to enjoy himself. His wife was not always out of humour, nor his home always uncomfortable; and in his breed of horses and dogs, and in sporting of every kind, he found no inconsiderable degree of domestic felicity.

For Marianne, however, in spite of his incivility in surviving her loss, he always retained that decided regard which interested him in every thing that befell her, and made her his secret standard of perfection in woman; and many a rising beauty would be slighted by him in after-days as bearing no comparison with Mrs Brandon.

Mrs Dashwood was prudent enough to remain at the cottage, without attempting a removal to Delaford; and fortunately for Sir John and Mrs Jennings, when Marianne was taken from them, Margaret had reached an age highly suitable for dancing, and not very ineligible for being supposed to have a lover.

Between Barton and Delaford, there was that constant communication which strong family affection would naturally dictate; and among the merits and the happiness of Elinor and Marianne, let it not be ranked as the least considerable, that though sisters, and living almost within sight of each other, they could live without disagreement between themselves, or producing coolness between their husbands.

\vfill
\begin{figure}[bh]
\centering
\includegraphics[width=.5\linewidth]{cameo}
\end{figure}

